\documentclass{leveldoc}
\DocUses{foundations,deflection_plates}
\DocFingerprint{ebada2d20f8008c6}

\begin{document}

\section{What the level asks}\label{sec:asks}

\leveldescription

One positive ion is launched along $+x$ from the left edge. A fixed positive
charge sits above the launch line. The player places up to three grounded
plates inside the marked region. The charge pushes the ion down; a grounded
plate between them carries the opposite induced charge and cuts that push on
the far side. There are no magnets.

The ion must enter the detector. The shot has no acceptance window, no kinetic
window and no gate. A flight ends when the ion touches the charge or a plate,
leaves the world, or enters the detector
(\xsecref{foundations}{sec:board}). A plate's obstacle is its box plus a
contact shell of thickness $\val{contact}$. The charge's obstacle radius is
$\val{Rcharge}$. The ion's radius is $0$.

The reference places one grounded plate between the charge and the launch
line. The ion drifts down and meets the detector's left side. With no plate
the charge drives the ion through the bottom edge of the world. This document
uses the deflection plates' rectangle collocation. The charge is external, on
the right-hand side, and there is no second plate, so there is no uniform gap.
The flight is integrated by the classical Runge--Kutta method.

\section{Setup}\label{sec:setup}

\setuptable

The speed of light is $c = \val{cLight}$. The ion has charge $q = \val{qPart}$,
mass $m = \val{mPart}$ and radius $0$, and it is launched with kinetic energy
$T_0 = \val{T0}$. The launch is at $(\val{xLaunch}, \val{yLaunch})$, straight
along $+x$. Then $\gamma_0 = \val{gamma0}$, $\abs{\vect{p}_0} = \val{p0}$ and
$v_0 = \val{v0} = \val{vOverC}\,c$ (\xeqref{foundations}{eq:launch}).

The fixed charge is $Q = \val{qFixed}$ at $(\val{xCharge}, \val{yCharge})$.
The level file does not list a plate size, so the plate is the engine's
default: length $\val{plateL}$ along $x$, thickness $\val{plateT}$ and height
$\val{plateH}$. The reference plate is centred at
$(\val{xPlateC}, \val{yPlateC})$, at potential $V = \val{vPlate}$, and it
occupies $\val{xPlate0} \le x \le \val{xPlate1}$,
$\val{yPlate0} \le y \le \val{yPlate1}$. That slab lies between the charge and
the launch line. The slider offers only this grounded potential.

The detector is $\val{xDetLo} \le x \le \val{xDetHi}$,
$\val{yDetLo} \le y \le \val{yDetHi}$. The world is $-1 \le x \le \val{xWorld}$,
$-1 \le y \le \val{yWorld}$. The reported time step is $h = \val{H}$. The mesh
uses rectangles of edge $\val{panelFine}$, $\val{nPanels}$ of them on the
upper half of the one plate.

\begin{figure}[ht]
\centering
\begin{tikzpicture}
\begin{axis}[axis equal image, width=0.92\linewidth,
  xmin=-1.5, xmax=32.5, ymin=-2, ymax=22.5,
  xlabel={$x$}, ylabel={$y$}, tick label style={font=\small}]
  \draw[black!40, dashed] (-1,-1) rectangle (\rawval{xWorld},\rawval{yWorld});
  \draw[black!25] (0,0) rectangle (30,20);
  \fill[green!50!black, fill opacity=0.18]
    (\rawval{xDetLo},\rawval{yDetLo}) rectangle (\rawval{xDetHi},\rawval{yDetHi});
  \draw[green!40!black]
    (\rawval{xDetLo},\rawval{yDetLo}) rectangle (\rawval{xDetHi},\rawval{yDetHi});
  \fill[black!20]
    (\rawval{xPlate0},\rawval{yPlate0}) rectangle (\rawval{xPlate1},\rawval{yPlate1});
  \draw[black!70]
    (\rawval{xPlate0},\rawval{yPlate0}) rectangle (\rawval{xPlate1},\rawval{yPlate1});
  \fill[red!70!black] (\rawval{xCharge},\rawval{yCharge}) circle[radius=\rawval{Rcharge}];
  \addplot[orange!80!black] table {pathbare.dat};
  \addplot[very thick, blue!60!black] table {pathref.dat};
  \draw[-{Stealth[length=5pt]}, thick] (\rawval{xLaunch},\rawval{yLaunch}) -- (\rawval{arrowX},\rawval{arrowY});
  \fill (\rawval{xLaunch},\rawval{yLaunch}) circle[radius=0.08];
\end{axis}
\end{tikzpicture}
\caption{The plate, the charge and the two flights. The red disc is the fixed
charge. The gray rectangle is the grounded plate, between the charge and the
launch line. The green box is the detector. Blue is the reference flight: the
ion drifts down and meets the left side of the detector. Orange is the same
launch with no plate: the charge pushes the ion down through the bottom edge
of the world, short of the detector. The blue path stays outside the plate.}
\label{fig:board}
\end{figure}

\section{Principles and model}\label{sec:model}

\begin{principles}
  \item The relativistic equation of motion \xeqref{foundations}{eq:eom}, the
    launch \xeqref{foundations}{eq:launch}, and energy in a static field
    \xeqref{foundations}{eq:energy}. There is no magnetic field.
  \item A player plate is a metal box at the potential the player sets
    (\physref{2.7}). Its surface charge is the deflection plates' collocation
    \xeqref{deflection_plates}{eq:collocation}, with the rectangle integral
    \xeqref{deflection_plates}{eq:rectangle}. The fixed charge is external.
  \item With the plate omitted the field is the point-charge field. That
    flight does not use a mesh.
  \item The board's rules (\xsecref{foundations}{sec:board}) and the classical
    Runge--Kutta method (\xsecref{foundations}{sec:numerics}).
  \item The radiated power \xeqref{foundations}{eq:larmor}, summed by Simpson's
    rule \xeqref{foundations}{eq:simpson}. It is not put back into the force.
\end{principles}

The reference has one grounded plate and the fixed charge. The state obeys
\xeqref{foundations}{eq:eom} with $\vect{B} = \vect{0}$ until the first ending
event.

\begin{assumption}{A static plate}\label{as:static}
The potential equals the grounded value at once, and the charge is fixed. The
launch speed is $\val{vOverC}\,c$. The game uses the same static plate
(\physref{2.7}).
\end{assumption}

\begin{assumption}{No radiation reaction}\label{as:radiation}
The radiated energy is not taken from the flight. On the reference flight it
is a fraction $\val{RefRadFrac}$ of $T_0$.
\end{assumption}

\begin{assumption}{No image of the ion}\label{as:image}
The ion's own surface charge is left out, as on the deflection plates and as
in the game (\physref{2.7}). The bound $q^2/d^2$, with $d$ floored at
$10^{-3}$ and taken relative to $\max(\abs{q\vect{E}}, T_0)$, reaches
$\val{RefImage}$ on the reference samples. The flight with no plate has no
metal, so its ratio is $\val{BareImage}$. The engine asks for the ratio to
stay below $10^{-10}$.
\end{assumption}

\begin{assumption}{Equal rectangles}\label{as:panels}
The mesh is equal rectangles of edge $\val{panelFine}$. Replacing the edge by
$\val{panelCoarse}$ moves an end state by $\val{errMesh}$. The flight with no
plate has no panels, so it does not depend on the mesh.
\end{assumption}

\begin{assumption}{The fixed step}\label{as:step}
The reported step is $h = \val{H}$. Halving it estimates an end state by
$\val{errStep}$. The tolerance on a state is ten times the larger of that
estimate and the panel-size estimate, and at least $10^{-9}$. Here the
panel-size estimate is the larger one, so the tolerance is $\val{tolFlight}$.
Between samples the ion moves at most $\val{stepReach}$. The clearance of the
plate is sampled at the ends of the steps and at the ending point, and it
reaches $\val{RefDmin}$.
\end{assumption}

\section{Derivation}\label{sec:derivation}

\subsection{The screen}\label{sec:plate}

A player plate is an electrode. The rectangles, the mirror and the collocation
are \xeqref{deflection_plates}{eq:rectangle} and
\xeqref{deflection_plates}{eq:collocation}. The plate potential on the
right-hand side is $\val{vPlate}$, and the fixed charge's potential is
subtracted there. The ion is not in the linear system. The reported mesh has
$\val{nPanels}$ upper rectangles. The residual at the centroids is
$\val{solveRes}$. On the face toward the beam, off the centroids, the
potential differs from $V$ by at most $\val{surfRes}$.

There is no second plate, so the uniform-gap field does not apply. The check
is the field on the launch line under the plate's centre, a distance
$\val{rSample}$ from the charge. The point-charge field there is
$\val{eCoulomb}$, along $-y$, and the potential is $\val{phiCoulomb}$. The
flight with no plate reproduces that field: the relative difference is
$\val{coulombRel}$. With the plate, the field at the same station is
$\val{eUnder}$, along $+y$, and the potential is $\val{phiUnder}$. The ratio
of the magnitudes is $\val{screenRel}$. The induced charge of the plate and
its mirror is $\val{qNet}$, opposite the fixed charge. The reversed field is
local. The march uses the panels, because the ion passes $\val{RefDmin}$ from
the metal, and the path's highest point is still the launch.

\subsection{The flight}\label{sec:flight}

The diagnostic is \xeqref{deflection_plates}{eq:flight-energy}. The state is
advanced by \xeqref{foundations}{eq:rk4} with step $h = \val{H}$. Events are
located as on the deflection plates: thirty-two Hermite samples, then sixty
bisections, obstacles before the world, and the world before the detector. The
step-halving estimate divides the change by $15$; the panel-size estimate is
the change itself when the edge doubles.

\section{The reference solution by hand}\label{sec:hand}

\begin{procedure}
  \item\label{st:launch} \emph{The launch} \xeqref{foundations}{eq:launch}.
    $\gamma_0 = \val{gamma0}$, $\abs{\vect{p}_0} = \val{p0}$, so
    $\vect{p}_0 = (\val{p0}, 0)$ and $v_0 = \val{v0}$.
  \item\label{st:plate} \emph{The screen.} The charge $Q = \val{qFixed}$ is at
    $(\val{xCharge}, \val{yCharge})$. The plate is centred at
    $(\val{xPlateC}, \val{yPlateC})$, at potential $\val{vPlate}$, with size
    $\val{plateL}$ by $\val{plateT}$ by $\val{plateH}$. Solve
    \xeqref{deflection_plates}{eq:collocation}. The centroid residual is
    $\val{solveRes}$ and the off-centroid discrepancy is $\val{surfRes}$. The
    induced charge is $\val{qNet}$. Under the plate the field is
    $\val{eUnder}$, the fraction $\val{screenRel}$ of the point-charge field,
    and of the opposite sign.
  \item\label{st:march} \emph{The march.} Advance
    \xeqref{deflection_plates}{eq:flight-energy} and
    \xeqref{foundations}{eq:eom} by \xeqref{foundations}{eq:rk4} with step
    $\val{H}$. The state every five time units:
    \begin{center}\small\val{states}\end{center}
  \item\label{st:end} \emph{The arrival.} The ion meets the \val{RefFace} face
    of the detector at $(\val{RefX}, \val{RefY})$, at $t = \val{RefT}$, with
    momentum $(\val{RefPx}, \val{RefPy})$, at an angle $\val{RefAng}$ degrees
    from $+x$. The abscissa is the detector's left edge, and the ordinate lies
    between $\val{yDetLo}$ and $\val{yDetHi}$. The run takes $\val{RefSteps}$
    steps. The path only falls: its highest point is the launch, and its
    lowest is $\val{RefYmin}$.
  \item\label{st:check} \emph{The checks.} $W$ drifts by at most
    $\val{RefWdev}$. The step-halving estimate is $\val{RefErr}$, and the
    larger of the two flights is $\val{errStep}$. The panel-size estimate is
    $\val{errMesh}$. The sampled clearance is $\val{RefDmin}$, outside the
    shell $\val{contact}$. The image-force ratio reaches $\val{RefImage}$.
  \item\label{st:rad} \emph{The radiation.} Simpson's rule on steps of $h/2$
    gives $\val{RefRad}$, a fraction $\val{RefRadFrac}$ of $T_0$. Halving that
    step again changes the integral by the relative amount $\val{RefRadErr}$.
    This integral is not a row of the comparison.
\end{procedure}

\begin{result}[The reference flight]
The ion enters the detector through its left side at
$(\val{RefX}, \val{RefY})$ at $t = \val{RefT}$, with momentum
$(\val{RefPx}, \val{RefPy})$. The level is solved.
\end{result}

\section{Comparison with the game}\label{sec:comparison}

\comparisontable

Eight comparisons are in the table: the reference outcome, time and end
position, and the outcome, time and end position with no plate. The tolerance
on a state is $\val{tolFlight}$. Ten times the panel-size estimate
$\val{errMesh}$ lies above $10^{-9}$, and it is larger than the step-halving
estimate $\val{errStep}$, so the tolerance is that product. The tolerance on a
state was fixed from these estimates before the comparison.

The largest absolute difference is the reference ordinate,
$\val{cmpRefYDiff}$. It lies above the panel-size estimate $\val{errMesh}$ and
under the tolerance. The reference time differs by $\val{cmpRefTDiff}$. The
bare time differs by $\val{cmpBareTDiff}$, the bare abscissa by
$\val{cmpBareXDiff}$ and the bare ordinate by $\val{cmpBareYDiff}$. The bare
field is the point charge, so the panel estimate does not apply to that
flight; those three differences lie under the tolerance. The outcomes agree.

The reference abscissa agrees exactly. It is the detector's left edge.

\section{Other solutions}\label{sec:others}

No plate, which is the empty placement in \code{placements.json}, does not
solve the level. The field is the point charge, and the ion is pushed along
$-y$. It leaves through the \val{BareFace} edge of the world at
$(\val{BareX}, \val{BareY})$, at $t = \val{BareT}$, with momentum
$(\val{BarePx}, \val{BarePy})$, at an angle $\val{BareAng}$ degrees from $+x$.
The ordinate $\val{BareY}$ is the bottom edge. The abscissa is short of the
detector, which begins at $\val{xDetLo}$. The energy drifts by $\val{BareWdev}$,
and the run takes $\val{BareSteps}$ steps.

The reference works because the grounded plate sits between the charge and the
launch. The ion still falls, from $\val{yLaunch}$ to $\val{RefY}$, and that
ordinate is inside the detector. This document integrates the reference plate
and the charge alone, not a second or third plate. The slider has no other
potential.

\section{What the game computes}\label{sec:engine}

A player plate is the same boundary-element electrode as a level plate, with
the default size when the level file omits it (\physref{2.7}). This document
collocates \xeqref{deflection_plates}{eq:collocation} on equal rectangles, and
the fixed charge enters the right-hand side. The end states move by
$\val{errMesh}$ when that edge doubles, and the reference ordinate in
\S\ref{sec:comparison} lies above that estimate and under ten times it. The
flight with no plate is the point-charge field: there is nothing to mesh, and
our sample agrees exactly with that Coulomb field.

The ion's image force is omitted on both sides. The bound sampled on the
reference flight reaches $\val{RefImage}$, under the engine's limit
$10^{-10}$. The flight in the engine is the adaptive method of \physref{5} and
\physref{6}, verified as in \physref{7}. Our step estimate is $\val{errStep}$.
Radiation is integrated and not applied to the momentum
(\xsecref{foundations}{sec:radiation}, \physref{8}).

\section{Further reading}\label{sec:reading}

The rectangle integral and the collocation are derived for the deflection
plates (\xeqref{deflection_plates}{eq:rectangle},
\xeqref{deflection_plates}{eq:collocation}). The motion and the Runge--Kutta
step are in the Foundations (\xsecref{foundations}{sec:numerics}). The
engine's player plates are \physref{2.7}.

\levelfileappendix

\end{document}
